\scanpage{129}
ment définir tout de suite
\[
\tag{206}
\left\{\begin{aligned}
\mathcal D_*^*&=\coprod_{d_0\geq d_1\cdots\geq d_r\geq0}D^*_{d_0,d_1,\ldots,d_r}\\
&=\text{image inverse de }D_*^*=\coprod_{d\in\mathbf N}D_d^*\text{ par}\\
&\quad\text{le morphisme }\unsure{0|o}:\mathcal D_*\longrightarrow D_*\\
&\quad\text{défini par les }D_{d_0,\ldots,d_r}\longrightarrow D_{d_0}\\
&\quad\text{i.e. défini par les }\mathcal D_r\longrightarrow D_*
\end{aligned}\right.
\]
\inserted{applications « premier sommet »}
(transposé de
\diagram{\begin{tikzcd}[ampersand replacement=\&, column sep=small,row sep=small]
\&\bullet\,0\arrow[d]\\
\underset{0}{\bullet}\arrow[ur]\&1\arrow[d]\\
\&2\arrow[d,no head]\\
\&\vdots\arrow[d]\\
\&r
\end{tikzcd}}
)

Nous allons définir
\[
\tag{207}
\left\{\begin{aligned}
P_{d_0,d_1,\ldots,d_r}&=\text{image inverse du }\overbrace{P_{d_0}}^{G_{d_0}\text{-torseur}}\text{ sur }D_{d_0}\\
&\qquad\text{par }D_{d_0,\ldots,d_r}\unsure{\text{---}|\longrightarrow}D_{d_0}\\
P^*_{d_0,d_1,\ldots,d_r}&=P_{d_0,\ldots,d_r}\mid D^*_{d_0,\ldots,d_r}
\end{aligned}\right.
\]
\unsure{c'est}\inserted{\unsure{là|sont}} des torseurs sous les Groupes analytiques
(tores)
\[
\tag{208}
\begin{aligned}
G_{d_0,\ldots,d_r}&=\text{image inverse de }G_{d_0}\\
&\quad\text{sur }\unsure{D_{d_0}|d_0}\text{, par }D_{d_0,\ldots,d_r}\unsure{\text{---}|\longrightarrow}D_{d_0}\\
G^*_{d_0,\ldots,d_r}&=G_{d_0,\ldots,d_r}\mid D^*_{d_0,\ldots,d_r}.
\end{aligned}
\]
On a donc un diagramme $2$-cartésien de groupoïdes
\[
(209)
\]
\diagram{\begin{tikzcd}[ampersand replacement=\&, column sep=large,row sep=large]
\bigl(\widetilde\Pi_{d_0- d_r}\overset{\mathrm{déf}}{=}\Pi_1P^*_{d_0,\ldots,d_r}
\arrow[r]\arrow[d]
\&\Pi_1P^*_{d_0}\bigl(\overset{\mathrm{déf}}{=}\widetilde\Pi_{d_0}\bigr)\arrow[d]\\
\Pi_{d_0,\ldots,d_r}=\Pi_1\unsure{D_{d_0,\ldots,d_r}}\arrow[r]
\&\Pi D_{d_0}\bigl(\overset{\mathrm{déf}}{=}\Pi_{d_0}\bigr)
\arrow[ul,phantom,"2\text{-cart.}" description]
\end{tikzcd}}
\margin{128}
